%% =====================================================================
%%  B3-T-21-05 -- what B3 contributes to "Strike the Shepherd"
%%  -------------------------------------------------------------------
%%  LAYER A: APPEND-ONLY. Written once at the entry's write, then left
%%  alone. If the source has more to say later, add a bullet -- do not
%%  rewrite. Material found ONLY here goes in the `unique` slot with
%%  @unique set to yes, which makes it undeletable (rule R10).
%% =====================================================================
%% @kind: source
%% @id: B3-T-21-05
%% @book: B3
%% @topic: T-21-05
%% @locator: Law 42, pp. 357--366
%% @status: distilled
%% @unique: no
%% @integrated: yes
%% @updated: 2026-09-19

\dossier{B3}{T-21-05}{\bookshort{B3}}{Law 42, pp. 357--366}

\begin{distilled}
  \item Trouble can often be traced to a single strong individual --- the stirrer, the malicious undercurrent of goodwill.
  \item Do not wait for the troubles they cause to multiply; do not try to negotiate --- they are untreatable, and negotiation mints more grievance.
  \item Neutralize their influence by isolating or banishing them; strike at the source and the sheep will scatter.
  \item Athens' ostracism institutionalised the removal of the one man whose ambition the city could not otherwise survive.
  \item Reversal: strike the wrong target, or the right one into martyrdom, and you have created the author the room lacked.
\end{distilled}
